\documentclass[11pt]{article}

\title{Math 603: Homework 3}
\author{Swayam Chube}
\date{Last Updated: \today}

\usepackage[utf8]{inputenc} % allow utf-8 input
\usepackage[T1]{fontenc}    % use 8-bit T1 fonts
\usepackage{hyperref}       % hyperlinks
\usepackage{url}            % simple URL typesetting
\usepackage{booktabs}       % professional-quality tables
\usepackage{amsfonts}       % blackboard math symbols
\usepackage{nicefrac}       % compact symbols for 1/2, etc.
\usepackage{microtype}      % microtypography
\usepackage{graphicx}
\usepackage{natbib}
\usepackage{doi}
\usepackage{amssymb}
\usepackage{bbm}
\usepackage{amsthm}
\usepackage{amsmath}
\usepackage{xcolor}
\usepackage{theoremref}
\usepackage{enumitem}
\usepackage{fouriernc}
\usepackage{mdframed}
\usepackage{mathrsfs}
\setlength{\marginparwidth}{2cm}
\usepackage{todonotes}
\usepackage{stmaryrd}
\usepackage[all,cmtip]{xy} % For diagrams, praise the Freyd-Mitchell theorem 
\usepackage{marvosym}
\usepackage{geometry}
\usepackage{titlesec}
\usepackage{mathtools}
\usepackage{tikz}
\usetikzlibrary{cd}
\usepackage{epigraph}
\setlength\epigraphwidth{0.4\textwidth}

\renewcommand{\qedsymbol}{$\blacksquare$}
% \renewcommand{\familydefault}{\sfdefault} % Do you want this font? 

% Uncomment to override  the `A preprint' in the header
% \renewcommand{\headeright}{}
% \renewcommand{\undertitle}{}
% \renewcommand{\shorttitle}{}

\hypersetup{
    pdfauthor={Swayam Chube},
    colorlinks=true,
	citecolor=blue,
}

\newtheoremstyle{thmstyle}%               % Name
  {}%                                     % Space above
  {}%                                     % Space below
  {}%                             % Body font
  {}%                                     % Indent amount
  {\bfseries\scshape}%                            % Theorem head font
  {.}%                                    % Punctuation after theorem head
  { }%                                    % Space after theorem head, ' ', or \newline
  {\thmname{#1}\thmnumber{ #2}\thmnote{ (#3)}}%                                     % Theorem head spec (can be left empty, meaning `normal')

\newtheoremstyle{defstyle}%               % Name
  {}%                                     % Space above
  {}%                                     % Space below
  {}%                                     % Body font
  {}%                                     % Indent amount
  {\bfseries\scshape}%                            % Theorem head font
  {.}%                                    % Punctuation after theorem head
  { }%                                    % Space after theorem head, ' ', or \newline
  {\thmname{#1}\thmnumber{ #2}\thmnote{ (#3)}}%                                     % Theorem head spec (can be left empty, meaning `normal')

\theoremstyle{thmstyle}
\newtheorem{theorem}{Theorem}[section]
\newtheorem{lemma}[theorem]{Lemma}
\newtheorem{proposition}[theorem]{Proposition}

\theoremstyle{defstyle}
\newtheorem{definition}[theorem]{Definition}
\newtheorem{corollary}[theorem]{Corollary}
\newtheorem{porism}[theorem]{Porism}
\newtheorem{remark}[theorem]{Remark}
\newtheorem{interlude}[theorem]{Interlude}
\newtheorem{example}[theorem]{Example}
\newtheorem*{notation}{Notation}
\newtheorem*{claim}{Claim}

% Common Algebraic Structures
\newcommand{\R}{\mathbb{R}}
\newcommand{\Q}{\mathbb{Q}}
\newcommand{\Z}{\mathbb{Z}}
\newcommand{\N}{\mathbb{N}}
\newcommand{\bbC}{\mathbb{C}} 
\newcommand{\K}{\mathbb{K}} % Base field which is either \R or \bbC
\newcommand{\F}{\mathbb{F}} % Base field which is either \R or \bbC
\newcommand{\calA}{\mathcal{A}} % Banach Algebras
\newcommand{\calB}{\mathcal{B}} % Banach Algebras
\newcommand{\calI}{\mathcal{I}} % ideal in a Banach algebra
\newcommand{\calJ}{\mathcal{J}} % ideal in a Banach algebra
\newcommand{\calR}{\mathcal{R}} % range of an operator
\newcommand{\frakM}{\mathfrak{M}} % sigma-algebra
\newcommand{\calO}{\mathcal{O}} % Ring of integers
\newcommand{\bbA}{\mathbb{A}} % Adele (or ring thereof)
\newcommand{\bbI}{\mathbb{I}} % Idele (or group thereof)

% Categories
\newcommand{\catTopp}{\mathbf{Top}_*}
\newcommand{\catGrp}{\mathbf{Grp}}
\newcommand{\catTopGrp}{\mathbf{TopGrp}}
\newcommand{\catSet}{\mathbf{Set}}
\newcommand{\catTop}{\mathbf{Top}}
\newcommand{\catRing}{\mathbf{Ring}}
\newcommand{\catCRing}{\mathbf{CRing}} % comm. rings
\newcommand{\catMod}{\mathbf{Mod}}
\newcommand{\catMon}{\mathbf{Mon}}
\newcommand{\catMan}{\mathbf{Man}} % manifolds
\newcommand{\catDiff}{\mathbf{Diff}} % smooth manifolds
\newcommand{\catAlg}{\mathbf{Alg}}
\newcommand{\catRep}{\mathbf{Rep}} % representations 
\newcommand{\catVec}{\mathbf{Vec}}

% Group and Representation Theory
\newcommand{\chr}{\operatorname{char}}
\newcommand{\Aut}{\operatorname{Aut}}
\newcommand{\GL}{\operatorname{GL}}
\newcommand{\im}{\operatorname{im}}
\newcommand{\tr}{\operatorname{tr}}
\newcommand{\id}{\mathbf{id}}
\newcommand{\cl}{\mathbf{cl}}
\newcommand{\Gal}{\operatorname{Gal}}
\newcommand{\Tr}{\operatorname{Tr}}
\newcommand{\sgn}{\operatorname{sgn}}
\newcommand{\Sym}{\operatorname{Sym}}
\newcommand{\Alt}{\operatorname{Alt}}

% Commutative and Homological Algebra
\newcommand{\spec}{\operatorname{spec}}
\newcommand{\mspec}{\operatorname{m-spec}}
\newcommand{\Spec}{\operatorname{Spec}}
\newcommand{\MaxSpec}{\operatorname{MaxSpec}}
\newcommand{\Tor}{\operatorname{Tor}}
\newcommand{\tor}{\operatorname{tor}}
\newcommand{\Ann}{\operatorname{Ann}}
\newcommand{\Supp}{\operatorname{Supp}}
\newcommand{\Hom}{\operatorname{Hom}}
\newcommand{\End}{\operatorname{End}}
\newcommand{\coker}{\operatorname{coker}}
\newcommand{\limit}{\varprojlim}
\newcommand{\colimit}{%
  \mathop{\mathpalette\colimit@{\rightarrowfill@\textstyle}}\nmlimits@
}
\makeatother


\newcommand{\fraka}{\mathfrak{a}} % ideal
\newcommand{\frakb}{\mathfrak{b}} % ideal
\newcommand{\frakc}{\mathfrak{c}} % ideal
\newcommand{\frakf}{\mathfrak{f}} % face map
\newcommand{\frakg}{\mathfrak{g}}
\newcommand{\frakh}{\mathfrak{h}}
\newcommand{\frakm}{\mathfrak{m}} % maximal ideal
\newcommand{\frakn}{\mathfrak{n}} % naximal ideal
\newcommand{\frakp}{\mathfrak{p}} % prime ideal
\newcommand{\frakq}{\mathfrak{q}} % qrime ideal
\newcommand{\fraks}{\mathfrak{s}}
\newcommand{\frakt}{\mathfrak{t}}
\newcommand{\frakz}{\mathfrak{z}}
\newcommand{\frakA}{\mathfrak{A}}
\newcommand{\frakB}{\mathfrak{B}}
\newcommand{\frakI}{\mathfrak{I}}
\newcommand{\frakJ}{\mathfrak{J}}
\newcommand{\frakK}{\mathfrak{K}}
\newcommand{\frakL}{\mathfrak{L}}
\newcommand{\frakN}{\mathfrak{N}} % nilradical 
\newcommand{\frakO}{\mathfrak{O}} % dedekind domain
\newcommand{\frakP}{\mathfrak{P}} % Prime ideal above
\newcommand{\frakQ}{\mathfrak{Q}} % Qrime ideal above 
\newcommand{\frakR}{\mathfrak{R}} % jacobson radical
\newcommand{\frakU}{\mathfrak{U}}
\newcommand{\frakV}{\mathfrak{V}}
\newcommand{\frakW}{\mathfrak{W}}
\newcommand{\frakX}{\mathfrak{X}}

% General/Differential/Algebraic Topology 
\newcommand{\scrA}{\mathscr{A}}
\newcommand{\scrB}{\mathscr{B}}
\newcommand{\scrF}{\mathscr{F}}
\newcommand{\scrM}{\mathscr{M}}
\newcommand{\scrN}{\mathscr{N}}
\newcommand{\scrP}{\mathscr{P}}
\newcommand{\scrO}{\mathscr{O}} % sheaf
\newcommand{\scrR}{\mathscr{R}}
\newcommand{\scrS}{\mathscr{S}}
\newcommand{\bbH}{\mathbb H}
\newcommand{\Int}{\operatorname{Int}}
\newcommand{\psimeq}{\simeq_p}
\newcommand{\wt}[1]{\widetilde{#1}}
\newcommand{\RP}{\mathbb{R}\text{P}}
\newcommand{\CP}{\mathbb{C}\text{P}}

% Miscellaneous
\newcommand{\wh}[1]{\widehat{#1}}
\newcommand{\calM}{\mathcal{M}}
\newcommand{\calP}{\mathcal{P}}
\newcommand{\onto}{\twoheadrightarrow}
\newcommand{\into}{\hookrightarrow}
\newcommand{\Gr}{\operatorname{Gr}}
\newcommand{\Span}{\operatorname{Span}}
\newcommand{\ev}{\operatorname{ev}}
% \newcommand{\weakto}{\stackrel{w}{\longrightarrow}}
\newcommand{\weakto}{\rightharpoonup}

\newcommand{\define}[1]{\textcolor{blue}{\textit{#1}}}
% \newcommand{\caution}[1]{\textcolor{red}{\textit{#1}}}
\newcommand{\important}[1]{\textcolor{red}{\textit{#1}}}
\renewcommand{\mod}{~\mathrm{mod}~}
\renewcommand{\le}{\leqslant}
\renewcommand{\leq}{\leqslant}
\renewcommand{\ge}{\geqslant}
\renewcommand{\geq}{\geqslant}
\newcommand{\Res}{\operatorname{Res}}
\newcommand{\floor}[1]{\left\lfloor #1\right\rfloor}
\newcommand{\ceil}[1]{\left\lceil #1\right\rceil}
\newcommand{\gl}{\mathfrak{gl}}
\newcommand{\ad}{\operatorname{ad}}
\newcommand{\Stab}{\operatorname{Stab}}
\newcommand{\bfX}{\mathbf{X}}
\newcommand{\Ind}{\operatorname{Ind}}
\newcommand{\bfG}{\mathbf{G}}
\newcommand{\rank}{\operatorname{rank}}
\newcommand{\calo}{\mathcal{o}}
\newcommand{\frako}{\mathfrak{o}}
\newcommand{\Cl}{\operatorname{Cl}}

\newcommand{\idim}{\operatorname{idim}}
\newcommand{\pdim}{\operatorname{pdim}}
\newcommand{\Ext}{\operatorname{Ext}}
\newcommand{\co}{\operatorname{co}}
\newcommand{\bfO}{\mathbf{O}}
\newcommand{\bfF}{\mathbf{F}} % Fitting Subgroup
\newcommand{\Syl}{\operatorname{Syl}}
\newcommand{\nor}{\vartriangleleft}
\newcommand{\noreq}{\trianglelefteqslant}
\newcommand{\subnor}{\nor\!\nor}
\newcommand{\Soc}{\operatorname{Soc}}
\newcommand{\core}{\operatorname{core}}
\newcommand{\Sd}{\operatorname{Sd}}
\newcommand{\mesh}{\operatorname{mesh}}
\newcommand{\sminus}{\setminus}
\newcommand{\diam}{\operatorname{diam}}
\newcommand{\Ass}{\operatorname{Ass}}
\newcommand{\projdim}{\operatorname{proj~dim}}
\newcommand{\injdim}{\operatorname{inj~dim}}
\newcommand{\gldim}{\operatorname{gl~dim}}
\newcommand{\embdim}{\operatorname{emb~dim}}
\newcommand{\hght}{\operatorname{ht}}
\newcommand{\depth}{\operatorname{depth}}
\newcommand{\ul}[1]{\underline{#1}}
\newcommand{\type}{\operatorname{type}}
\newcommand{\CM}{\operatorname{CM}}
\newcommand{\cech}[1]{\mathbin{\check{#1}}}
\newcommand{\cdim}{\operatorname{cdim}}
\newcommand{\Der}{\operatorname{Der}}
\newcommand{\trdeg}{\operatorname{trdeg}}
\newcommand{\calE}{\mathcal{E}}
\newcommand{\calF}{\mathcal{F}}
\newcommand{\calT}{\mathcal{T}}
\newcommand{\calN}{\mathcal{N}}
\newcommand{\calL}{\mathcal{L}}
\renewcommand{\Re}{\operatorname{Re}}
\renewcommand{\Im}{\operatorname{Im}}
\newcommand{\gr}{\operatorname{gr}}
\newcommand{\dist}{\operatorname{dist}}
\newcommand{\calH}{\mathcal{H}}
\newcommand{\spr}{\operatorname{spr}}

\geometry {
    margin = 1in
}

\titleformat
{\section}
[block]
{\Large\bfseries\sffamily}
{\S\thesection}
{0.5em}
{\centering}
[]


\titleformat
{\subsection}
[block]
{\normalfont\bfseries\sffamily}
{\S\S}
{0.5em}
{\centering}
[]

\begin{document}
\maketitle 

\section{Problem 1}

\section{Problem 2}

\section{Problem 3}

Suppose first that $f$ is weakly analytic. Fix $z_0\in\Omega$. We shall show that the limit 
\begin{equation*}
    \lim_{z\to z_0}\frac{f(z) - f(z_0)}{z - z_0}
\end{equation*}
exists in $X$. Let $R > 0$ be such that the closed ball $\overline B_R(z_0)$ is contained in $\Omega$, and let $\Gamma\colon [0, 2\pi]\to\Omega$ denote the positively oriented circle $\Gamma(t) = z_0 + Re^{\iota t}$. Let $\varphi\in X^\ast$. Since $\varphi \circ f\colon\Omega\to\bbC$ is holomorphic, using Cauchy's Integral Formula, we have 
\begin{equation*}
    \varphi(f(z)) = \frac{1}{2\pi\iota}\oint_\Gamma \frac{\varphi(f(\zeta))}{\zeta - z}~\mathrm d\zeta
\end{equation*}
for every $z\in B_R(z_0)$, so that 
\begin{align*}
    \frac{\varphi(f(z)) - \varphi(f(z_0))}{z - z_0} &= \frac{1}{2\pi\iota(z - z_0)}\oint_\Gamma\frac{\varphi(f(\zeta))}{\zeta - z} - \frac{\varphi(f(\zeta))}{\zeta - z_0}~\mathrm d\zeta\\
    &= \frac{1}{2\pi\iota(z - z_0)}\oint_\Gamma\varphi(f(\zeta))\frac{z - z_0}{(\zeta - z)(\zeta - z_0)}~\mathrm d\zeta\\
    &= \frac{1}{2\pi\iota}\oint_\Gamma\frac{\varphi(f(\zeta))}{(\zeta - z)(\zeta - z_0)}~\mathrm d\zeta.
\end{align*}
Therefore, for $z, z'\in B_{\frac{R}{2}}(z_0)$, we have 
\begin{align*}
    \frac{\varphi(f(z')) - \varphi(f(z_0))}{z' - z_0} - \frac{\varphi(f(z)) - \varphi(f(z_0))}{z - z_0} &= \frac{1}{2\pi\iota}\oint_{\Gamma}\frac{\varphi(f(\zeta))}{(\zeta - z')(\zeta - z_0)} - \frac{\varphi(f(\zeta))}{(\zeta - z)(\zeta - z_0)}~\mathrm d\zeta \\
    &= \frac{1}{2\pi\iota}\int_\Gamma\frac{\varphi(f(\zeta))\left(\zeta - z - (\zeta - z')\right)}{(\zeta - z')(\zeta - z_0)(\zeta - z)}~\mathrm d\zeta\\
    &= \frac{1}{2\pi\iota}\oint_\Gamma\frac{\varphi(f(\zeta))(z' - z)}{(\zeta - z')(\zeta - z)(\zeta - z_0)}~\mathrm d\zeta.
\end{align*}
Rewriting the left hand side,
\begin{equation}\label{cauchy-esque}\tag{$\dagger$}
    \varphi\left(\frac{f(z') - f(z_0)}{z' - z_0} - \frac{f(z) - f(z_0)}{z - z_0}\right) = \frac{1}{2\pi\iota}\oint_\Gamma\frac{\varphi(f(\zeta))(z' - z)}{(\zeta - z')(\zeta - z)(\zeta - z_0)}~\mathrm d\zeta
\end{equation}
for all $z, z'\in B_{\frac{R}{2}}(z_0)$.

\begin{claim}
    The set $\left\{f(\zeta)\colon\zeta\in\Gamma\right\}$ is bounded.
\end{claim}
\begin{proof}
    We shall show this using the Banach-Steinhaus Theorem (i.e., the Uniform Boundedness Principle). Indeed, for every $\varphi\in X^\ast$, the function $\varphi\circ f$ is continuous and since $\Gamma$ is compact, there is a constant $B_\varphi > 0$ such that $|\varphi(f(\zeta))|\le B_\varphi$ for every $\zeta\in\Gamma$. Let $\Phi\colon X\to X^{\ast\ast}$ denote the canonical isometric map (called the ``evaluation map'' sometimes). Note that $\varphi(f(\zeta)) = \Phi(f(\zeta))(\varphi)$. We have shown that the bounded linear functionals $\Phi(f(\zeta))\in X^{\ast\ast}$ for $\zeta\in\Gamma$ are pointwise bounded on $X^\ast$, and hence, due to the Banach-Steinhaus Theorem, there is a constant $C > 0$ such that $\|\Phi(f(\zeta))\|\le C$ for every $\zeta\in\Gamma$. Using the fact that $\Phi$ is an isometry, we see that $\|f(\zeta)\|\le C$ for every $\zeta\in\Gamma$, as desired.
\end{proof}

Using the above in \eqref{cauchy-esque}, we have 
\begin{align*}
    \left|\varphi\left(\frac{f(z') - f(z_0)}{z' - z_0} - \frac{f(z) - f(z_0)}{z - z_0}\right)\right| &= \frac{1}{2\pi}\left|\oint_{\Gamma}\frac{\varphi(f(\zeta))(z' - z)}{(\zeta - z')(\zeta - z)(\zeta - z_0)}~\mathrm d\zeta\right|\\
    &= \frac{1}{2\pi}\oint_\Gamma\frac{|\varphi(f(\zeta))| |z' - z|}{|\zeta - z'||\zeta - z||\zeta - z_0|}~\mathrm d|\zeta|\\
    &\le \frac{1}{2\pi}\oint_\Gamma\frac{\|\varphi\|\|f(\zeta)\| |z' - z|}{|\zeta - z'||\zeta - z||\zeta - z_0|}~\mathrm d|\zeta|\\
    &\le\frac{1}{2\pi}\oint_\Gamma\frac{\|\varphi\|\cdot C\cdot |z' - z|}{\frac{R^3}{4}}~\mathrm d|\zeta|\\
    &=\frac{4C\|\varphi\|}{R^2}|z' - z|,
\end{align*}
for all $\varphi\in X^\ast$.
Let
\begin{equation*}
    x = \frac{f(z') - f(z_0)}{z' - z_0} - \frac{f(z) - f(z_0)}{z - z_0}\in X.
\end{equation*}
Then the above can be rewritten as: 
\begin{equation*}
    \left|\Phi(x)(\varphi)\right|\le\frac{4C|z' - z|}{R^2}\|\varphi\|
\end{equation*}
for all $\varphi\in X^\ast$. It follows that 
\begin{equation*}
    \left\|\frac{f(z') - f(z_0)}{z' - z_0} - \frac{f(z) - f(z_0)}{z - z_0}\right\| = \|x\| = \|\Phi(x)\|\le\frac{4C}{R^2}|z' - z|,
\end{equation*}
for all $z, z'\in B_{\frac{R}{2}}(z_0)$. Let $g\colon \Omega\setminus\{z_0\}\to X$ denote the mapping 
\begin{equation*}
    g(z) = \frac{f(z) - f(z_0)}{z - z_0}.
\end{equation*}
For any sequence $(z_n)_{n = 1}^\infty$ in $B_{\frac{R}{2}}(z_0)\setminus\{z_0\}$, we have 
\begin{equation*}
    \|g(z_n) - g(z_m)\|\le\frac{4C}{R^2}|z_n - z_m|,
\end{equation*}
so that the sequence $\left(g(z_n)\right)_{n = 1}^\infty$ is Cauchy, and hence, converges since $X$ is complete. Since this holds for every sequence $(z_n)_{n = 1}^\infty$ converging to $z_0$ (because any such sequence will eventually be contained in the $\frac R2$-ball centered at $z_0$), it follows that the limit 
\begin{equation*}
    \lim_{z\to z_0}\frac{f(z) - f(z_0)}{z - z_0}
\end{equation*}
exists. In the above analytis $z_0\in\Omega$ was arbitrary, and hence, $f$ is strongly analytic.

Conversely, suppose $f$ is strongly analytic. Then for every $\varphi\in X^\ast$ and $z_0\in\Omega$, we have that for every $z\in\Omega$ with $z\ne z_0$, 
\begin{equation*}
    \lim_{z\to z_0}\frac{\varphi(f(z)) - \varphi(f(z_0))}{z - z_0} = \lim_{z\to z_0}\varphi\left(\frac{f(z) - f(z_0)}{z - z_0}\right) = \varphi\left(\lim_{z\to z_0}\frac{f(z) - f(z_0)}{z - z_0}\right) = \varphi(f'(z_0)),
\end{equation*}
where we may swap the limit with $\varphi$ because it is continuous. Further, $f'(z_0)$ is a shorthand for the limit $\lim_{z\to z_0}\frac{f(z) - f(z_0)}{z - z_0}$. This shows that $f$ is weakly analytic, as desired.

\section{Problem 4}

Let $\mathbf K$ denote the ``integral operator'', that is, 
\begin{equation*}
    \left(\mathbf K f\right)(t) = \int_0^t f(s)~\mathrm ds,
\end{equation*}
which is obtained by setting the kernel $K(s, t)$ to be identically $1$. We first prove the result for $\mathbf K$. To this end, we shall show that the spectral radius of $\mathbf K$ is zero. Using the spectral radius formula,
\begin{equation*}
    \spr(\mathbf K) = \lim_{n\to\infty}\|\mathbf K^n\|^\frac{1}{n}
\end{equation*}
We shall show by induction on $n\ge 1$ that 
\begin{equation*}
    \left(\mathbf K^n f\right)(t)  = \int_0^t \frac{(t - s)^{n - 1}}{(n - 1)!}f(s)~\mathrm ds.
\end{equation*}
The base case with $n = 1$ is clear. Suppose now that the claim holds for $n\ge 1$. Then for $f\in C([0, 1])$, 
\begin{equation*}
    \left(\mathbf K^{n + 1}f\right)(t) = \int_0^t\left(\mathbf K^n f\right)(s)~\mathrm ds = \int_0^t\int_0^s\frac{(s - u)^{n - 1}}{(n - 1)!}f(u)~\mathrm du~\mathrm ds.
\end{equation*}
By Fubini's Theorem, the rightmost integral is equal to 
\begin{equation*}
    \iint_E\frac{(s - u)^{n - 1}}{(n - 1)!}f(u)~\mathrm du~\mathrm ds,
\end{equation*}
where $E = \left\{(u, s)\in\R^2\colon 0\le u\le s\le t\right\}$. Using Fubini's Theorem again, the integral can be rewritten as 
\begin{equation*}
    \int_{0}^t\int_u^t\frac{(s - u)^{n - 1}}{(n - 1)!}f(u)~\mathrm ds~\mathrm du = \int_0^t\frac{(t - u)^n}{n!}f(u)~\mathrm du,
\end{equation*}
as desired. 

Hence, for $f\in C([0, 1])$, we have 
\begin{equation*}
    \left|(\mathbf K^n f)(t)\right|\le\int_0^t\frac{(t - s)^{n - 1}}{(n - 1)!}|f(s)|~\mathrm ds\le\int_0^1\frac{1}{(n - 1)!}\|f\|~\mathrm ds = \frac{\|f\|}{(n - 1)!},
\end{equation*}
so that 
\begin{equation*}
    \|\mathbf K^n f\|\le\frac{\|f\|}{(n - 1)!}\implies\|\mathbf K^n\|\le\frac{1}{(n - 1)!}.
\end{equation*}
This shows that 
\begin{equation*}
    \lim_{n\to\infty}\|\mathbf K^n\|^{\frac{1}{n}}\le\lim_{n\to\infty}\frac{1}{\left((n - 1)!\right)^{\frac{1}{n}}} = 0.
\end{equation*}
With this in hand, we turn our attention to more general continuous kernels. Let $K\in C\left([0, 1]\times [0, 1]\right)$ be a continuous function. Since $K$ is continuous on a compact set, there is a constant $C > 0$ such that $|K(s, t)|\le C$ for all $s, t\in [0, 1]$. Let $V\colon C([0, 1])\to C([0, 1])$ denote the operator 
\begin{equation*}
    (Vf)(t) = \int_0^t K(t, s)f(s)~\mathrm ds.
\end{equation*}
Then for every $f\in C([0, 1])$ and $n\ge 1$, we have 
\begin{equation*}
    (V^nf)(t) = \int_0^t K(t, s_{n - 1})\int_0^{s_{n - 1}}K(s_{n - 1}, s_{n - 2})\int_0^{s_{n - 2}}\cdots\int_0^{s_1} K(s_1, s_0) f(s_0)~\mathrm ds_0\cdots\mathrm ds_{n - 1},
\end{equation*}
and hence, 
\begin{align*}
    \left|(V^n f)(t)\right|&\le\int_0^t |K(t, s_{n - 1})|\int_0^{s_{n - 1}} |K(s_{n - 1}, s_{n - 2})| \int_0^{s_{n - 2}}\cdots\int_0^{s_1} |K(s_1, s_0)||f(s_0)|~\mathrm ds_0\cdots\mathrm ds_{n - 1}\\
    &\le C^n\int_0^t \int_0^{s_{n - 1}}\int_0^{s_{n - 2}}\cdots\int_0^{s_1}|f(s_0)|~\mathrm ds_0\cdots\mathrm ds_{n - 1}\\
    &= C^n\left(\mathbf K^n|f|\right)(t)\\
    &\le C^n\|\mathbf K^n f\|\le C^n\|\mathbf K^n\|\|f\|.
\end{align*}
Thus $\|V^n f\|\le C^n\|\mathbf K^n\|\|f\|$. Since this inequality holds for all $f\in C([0, 1])$, it follows that $\|V^n\|\le C^n\|\mathbf K^n\|$, consequently, 
\begin{equation*}
    \spr(V) = \lim_{n\to\infty}\|V^n\|^{\frac{1}{n}} = \lim_{n\to\infty} C\|\mathbf K^n\|^{\frac{1}{n}} = 0.
\end{equation*}
Due to Gelfand-Mazur, we know that the spectrum of $V$ is non-empty, and hence, must be a singleton. All that remains to be shown is that $0\in\sigma(V)$. To this end, it suffices to show that $V$ is a compact operator. Indeed, let $(f_n)_{n = 1}^\infty$ be a bounded sequence in $C([0, 1])$. We shall show that the sequence $\left(Vf_n\right)_{n = 1}^\infty$ admits a uniformly convergent subsequence by invoking the Arzel\`a-Ascoli Theorem, for which we must show that $(Vf_n)_{n = 1}^\infty$ is pointwise bounded and equicontinuous. 
Suppose $B > 0$ is such that $\|f_n\|\le B$ for all $n\ge 1$. Then for all $n\ge 1$ and $t\in [0, 1]$, we have 
\begin{equation*}
    |(Vf)(t)|\le\left|\int_0^t K(t, s)f(s)~\mathrm ds\right|\le \int_0^t |K(t, s)| |f(s)|~\mathrm ds\le \int_0^1 C\cdot B~\mathrm ds = C\cdot B.
\end{equation*}
This shows that the $Vf_n$'s are pointwise bounded. We shall show equicontinuity next. Let $\varepsilon > 0$. Since $K$ is continuous on a compact set $[0, 1]\times[0, 1]$, it is uniformly continuous, so that there is a $\delta > 0$ such that whenever $|(x, x') - (y, y')| < \delta$, we have $\left|K(x, x') - K(y, y')\right| < \frac{\varepsilon}{4B}$. Further, we may choose $\delta < \frac{\varepsilon}{4BC}$. Now, for $|x - y| < \delta$ and $n\ge 1$, we have 
\begin{align*}
    \left|(Vf_n)(x) - (Vf_n)(y)\right| &= \left|\int_0^y K(y, s)f(s)~\mathrm ds - \int_0^x K(x, s)f(s)~\mathrm ds\right|\\
    &= \left|\int_0^x K(y, s)f(s) - K(x, s)f(s)~\mathrm ds + \int_x^y K(y, s)f(s)~\mathrm ds\right|\\
    &\le \int_0^x |K(y, s) - K(x, s)||f(s)|~\mathrm ds + \int_x^y |K(y, s)||f(s)|~\mathrm ds\\
    &\le \int_0^1 |K(y, s) - K(x, s)||f(s)|~\mathrm ds + \int_x^y |K(y, s)||f(s)|~\mathrm ds\\
    &\le\int_0^1\frac{\varepsilon}{4B}\cdot B~\mathrm ds + |x - y|\cdot C\cdot B\\
    &\le \frac{\varepsilon}{2} < \varepsilon.
\end{align*}
This establishes equicontinuity, and hence, by the Arzel\`a-Ascoli Theorem, $\left(Vf_n\right)_{n\ge 1}$ admits a uniformly convergent subsequence, thereby showing that $V$ is compact. Since $0$ is in the spectrum of every compact operator, we have shown that $\sigma(V) = \{0\}$.

\section{Problem 5}

\subsection{Part (a)}

Suppose $0\ne\lambda\notin \sigma(AB)$, that is, $\lambda I - AB$ is invertible. Set 
\begin{equation*}
    T = \frac{1}{\lambda}\left(I + B\left(\lambda I - AB\right)^{-1}A\right).
\end{equation*}
Then 
\begin{align*}
    T(\lambda I - BA) &= \frac{1}{\lambda}\left(I + B(\lambda I - AB)^{-1}A\right)(\lambda I - BA)\\
    &= \frac{1}{\lambda}\left(\lambda I - BA + B(\lambda I - AB)^{-1}\left(\lambda A - ABA\right)\right)\\
    &=\frac{1}{\lambda}\left(\lambda I - BA + B(\lambda I - AB)^{-1}\left(\lambda I - AB\right)A\right)\\
    &= \frac{1}{\lambda}\left(\lambda I - BA + BA\right)\\
    &= I,
\end{align*}
and 
\begin{align*}
    (\lambda I - BA)T &= \frac{1}{\lambda}(I - BA)\left(I + B(\lambda I - AB)^{-1}A\right)\\
    &= \frac{1}{\lambda}\left(\lambda I - BA + (\lambda I - BA)B(\lambda I - AB)^{-1}A\right)\\
    &= \frac{1}{\lambda}\left(\lambda I - BA + (\lambda B - BAB)(\lambda I - AB)^{-1}A\right)\\
    &= \frac{1}{\lambda}\left(\lambda I - BA + B(\lambda I - AB)(\lambda I - AB)^{-1}A\right)\\
    &= \frac{1}{\lambda}\left(\lambda I - BA + BA\right)\\
    &= I.
\end{align*}
Thus $\lambda I - BA$ is invertible with inverse $T$. In other words, $\lambda\notin\sigma(BA)\setminus\{0\}$. Equivalently, we have shown that $\sigma(BA)\setminus\{0\}\subseteq\sigma(AB)\setminus\{0\}$. Interchanging the roles of $A$ and $B$ and repeating the above argument, we establish the reverse inclusion: $\sigma(AB)\setminus\{0\}\subseteq\sigma(BA)\setminus\{0\}$, whence we must have equality $\sigma(AB)\setminus\{0\} = \sigma(BA)\setminus\{0\}$.

\subsection{Part (b)}

Suppose $0\ne\lambda$ is an eigenvalue of $AB$, that is, there exists $0\ne x\in X$ such that $ABx = \lambda x$. Note that $Bx \ne 0$, else $\lambda x = 0$, which would force $\lambda = 0$, which is absurd. Set $y = Bx$. Then $BAy = BABx = \lambda Bx = \lambda y$, so that $\lambda$ is an eigenvalue of $BA$. Thus every non-zero eigenvalue of $AB$ is an eigenvalue of $BA$.

Interchanging the roles of $A$ and $B$, we see that every non-zero eigenvalue of $BA$ is also an eigenvalue of $AB$. Hence, $AB$ and $BA$ have the same non-zero eigenvalues.

\section{Problem 6} 

Let $X$ be a Banach space and $T\in\mathscr L(X)$ be a surjective bounded linear map. We shall show that there exists an $\varepsilon > 0$ such that whenever $S\in\mathscr L(X)$ with $\|T - S\| < \varepsilon$, $S$ is surjective. This statement is equivalent to the statement of the theorem with $A = T - S$. 

First, using the open mapping theorem, the image of the open unit ball of $X$ under $T$ is an open set containing $0$, and hence, there is an $r > 0$ such that the ball of radius $r$, $B_r(0)$ is contained in $T\left(B_1(0)\right)$. Now, given any $0\ne y\in X$, the element $\frac{r}{2\|y\|}y$ is contained in the ball $B_r(0)$, so that there exists $0\ne x\in B_1(0)$ with $Tx = \frac{r}{2\|y\|}y$. In particular, $T\left(\frac{2\|y\|}{r}x\right) = y$. Set $z = \frac{2\|y\|}{r}x$. In conclusion, we have shown that for every $y\in X$, there exists a $z\in X$ with $\|z\|\le\frac{2}{r}\|y\|$ such that $Tz = y$. Set $c = \frac{2}{r}$. We isolate our conclusion:
\begin{equation*}\label{conclusion}\tag{$\star$}
    \text{for every $y\in X$ there exists $x\in X$ with $\|x\|\le c\|y\|$ such that $Tx = y$}
\end{equation*}

Let $0 < \varepsilon < \frac{1}{2c}$ and let $S\in\mathscr L(X)$ with $\|T - S\| < \varepsilon$. We shall show that $S$ is surjective. Indeed, let $y\in X$ and set $y_0 = y$. Using \eqref{conclusion}, there exists $x_1\in X$ such that $Tx_1 = y_0$ and $\|x_1\|\le c\|y_0\| = c\|y\|$. Next, set $y_1 = y_0 - Sx_1$, so that 
\begin{equation*}
    \|y_1\| = \left\|Tx_1 - Sx_1\right\| = \left\|(T - S)(x_1)\right\|\le \varepsilon\|x_1\|\le c\varepsilon\|y\|.
\end{equation*}
Using \eqref{conclusion} again, there exists $x_2\in X$ with $Tx_2 = y_1$ and $\|x_2\|\le c\|y_1\|\le c^2\varepsilon \|y\|$. Set $y_2 = y_1 - Sx_2$, and using the same argument as above, we note that $\|y_2\|\le (c\varepsilon)^2 \|y\|$. Continuing this process \emph{ad infinitum}, we obtain two sequences $(x_n)_{n\ge 1}$ and $(y_n)_{n\ge 0}$ such that 
\begin{enumerate}[label=(\roman*)]
    \item $y_0 = y$,
    \item $\|x_n\|\le c^n\varepsilon^{n - 1}\|y\|$ and $\|y_n\|\le (c\varepsilon)^n \|y\|$ for all $n\ge 1$, and 
    \item $y_n = y_{n - 1} - Sx_n$ for all $n\ge 1$.
\end{enumerate}
Note that the sequence $(x_n)_{n\ge 1}$ is absolutely summable, since 
\begin{equation*}
    \sum_{n\ge 1}\|x_n\|\le \sum_{n\ge 1}c^n\varepsilon^{n - 1}\|y\| = c\sum_{n\ge 1}(c\varepsilon)^{n - 1} < \infty,
\end{equation*}
since $c\varepsilon < \frac{1}{2}$. Being absolutely summable, its partial sums 
\begin{equation*}
    z_N\coloneq\sum_{n = 1}^N x_n
\end{equation*}
must form a Cauchy sequence in $X$, owing to the inequality: 
\begin{equation*}
    \|z_N - z_M\|\le\sum_{n = M + 1}^N \|x_n\|.
\end{equation*}
Since $X$ is complete, the partial sums converge to some $x\in X$. We claim that $Sx = y$. Indeed, note that 
\begin{equation*}
    Sz_N = \sum_{n = 1}^N Sx_n = \sum_{n = 1}^N (y_{n - 1} - y_n) = y_0 - y_N = y - y_N.
\end{equation*}
But since $\|y_n\|\to 0$ as $n\to\infty$ (again using that $c\varepsilon < \frac{1}{2}$), we see that 
\begin{equation*}
    Sx = \lim_{N\to\infty} Sz_n = \lim_{N\to\infty} y - y_N = y,
\end{equation*}
and hence $S$ is surjective, as desired.
\end{document}